% kernel_redteam.tex -- round 495, PRIME.RDAGGER.KERNEL_LOEWNER_REDTEAM.01
% Adversarial review of r494.  Verdict CONFIRMED(c=2.1e-3).
% Companion: verify_kernel_redteam.py.
% Discovery: experiments/tfpt-discovery/kernel_redteam_probe.py.
% NO RH CLAIM.  NO anti-RH claim.
\documentclass[11pt]{article}

\usepackage[a4paper,margin=2.35cm]{geometry}
\usepackage{amsmath,amssymb,amsthm}
\usepackage{booktabs,array,microtype,xcolor}
\usepackage[T1]{fontenc}
\usepackage[colorlinks=true,linkcolor=blue!50!black,urlcolor=blue!50!black]{hyperref}

\newtheorem{lemma}{Lemma}
\newtheorem{theorem}{Theorem}
\theoremstyle{remark}
\newtheorem{remark}{Remark}
\newcommand{\QW}{Q_{\mathrm W}}
\newcommand{\CC}{\ensuremath{\mathrm{CC}}}
\setlength{\emergencystretch}{2em}

\title{\bfseries Adversarial audit of the kernel--Loewner floor\\[2mm]
\large PRIME.RDAGGER.KERNEL\_LOEWNER\_REDTEAM.01 --- round 495}
\author{Stefan Hamann}
\date{August 31, 2026}

\begin{document}
\maketitle

\begin{abstract}\noindent
This round tries to kill the r494 assertion
$\lambda_*(0.3)\ge2.1\cdot10^{-3}$ from first principles.
The identity was rederived from the defining digamma multiplier and
cross-checked on five new compactly supported functions at 50-plus
decimal precision.  Translation boundary strips were checked exactly
over $\mathbb Q$.  A false doubled-$c_L$ world is detected as negative.
An independent Clenshaw--Curtis Nystr\"om route reproduces the
$401$-section eigenvalue and Hilbert--Schmidt tail, including a
nonzero complementary block that could not be discarded.
No attack closes.  The outward r494 balance is
\[
 \begin{split}
 &3.69800804029-2.19240491114-1.5013673873\\
 &\qquad{}-3(7.0457565325\cdot10^{-4})
 =2.12201489024986\cdot10^{-3}.
 \end{split}
\]
\textbf{Verdict: \texttt{CONFIRMED($c=2.1000\mathrm e{-3}$)}.}
\end{abstract}

\begin{center}
\fbox{\parbox{0.94\linewidth}{%
\textbf{Claim boundary.}
This is an adversarial audit of one unconditional fixed-support Weil
inequality in the classical $2L<\log2$ zone.  It is not an RH claim,
not an anti-RH claim, and creates no ledger row.}}
\end{center}

\section{A1: independent derivation of the identity}

Use
\[
 \widehat h(t)=\int_{\mathbb R}h(u)e^{itu}\,du,\qquad
 g(x)=\int_{\mathbb R}h(u)h(u-x)\,du.
\]
Then $(2\pi)^{-1}\int|\widehat h(t)|^2\cos(tx)\,dt=g(x)$.
Starting from the Gauss representation, independently of the r494
note,
\[
 \psi(z)=-\gamma+\int_0^\infty
 \frac{e^{-v}-e^{-zv}}{1-e^{-v}}\,dv,
 \qquad \operatorname{Re}z>0,
\]
and putting $z=\tfrac14+\tfrac i2t$, $v=2x$, gives
\[
 \sigma_A(t)=-\gamma-\log\pi+
 \int_0^\infty
 \frac{2e^{-2x}-2e^{-x/2}\cos(tx)}{1-e^{-2x}}\,dx.
\]
The cosine coefficient is
\[
 k(x)=\frac{2e^{-x/2}}{1-e^{-2x}}
     =\frac{e^{x/2}}{\sinh x}.
\]
Plancherel therefore yields
\[
 A(h)=\int_0^\infty
 \left(\frac{e^{-2x}}xg(0)-k(x)g(x)\right)dx
 -(\log\pi)g(0).
\]
For $\operatorname{supp}h\subset[-L,L]$, $g(x)=0$ for $x>2L$.
The regularized elementary identity
\[
 \lim_{\varepsilon\downarrow0}
 \left(\int_\varepsilon^{2L}\frac{dx}{x}
 -\int_\varepsilon^\infty\frac{e^{-2x}}x\,dx\right)
 =\log(4L)+\gamma
\]
then gives
\[
 A(h)=\int_0^{2L}k(x)(g(0)-g(x))\,dx-c_Lg(0),
\]
\[
 c_L=\int_0^{2L}\left(k(x)-\frac1x\right)dx
 +\log(4L)+\gamma+\log\pi
 =2.1924049111319992552529\ldots .
\]

The pole sign is also forced rather than assumed.  Analytic
continuation gives
\[
 \widehat g(i/2)=\widehat h(i/2)\widehat h(-i/2)
 =(C-S)(C+S)=C^2-S^2,
\]
where $C=\langle h,\cosh(\cdot/2)\rangle$ and
$S=\langle h,\sinh(\cdot/2)\rangle$.  Thus
\[
 \Pi(h)=2(C^2-S^2).
\]
It is a product of analytic values, not
$|\widehat h(i/2)|^2$; the negative odd sign in r494 is correct.
Since $e^{2L}<2$ at $L=0.3$, there is no prime term.

\subsection*{Five non-calibration cross-tests}

Set $u=x/L$ and
$h_j(x)=(1-u^2)^8p_j(u)$ on $[-L,L]$, zero outside, with
\[
\begin{split}
p_1&=1,\quad p_2=u,\quad
p_3=1+\tfrac7{10}u-\tfrac25u^2,\\
p_4&=1-2u^2+\tfrac3{10}u^3,\quad
p_5=1-\tfrac13u+\tfrac15u^2-\tfrac17u^3.
\end{split}
\]
These are not the r494 calibration modes.  The frequency side uses
$\operatorname{Re}\psi(\tfrac14+it/2)-\log\pi$ directly at 60 working
digits; the $x$ side integrates the independently derived kernel.

\begin{center}
\small
\begin{tabular}{@{}rrrr@{}}
\toprule
$j$ & $\QW^{\rm freq}/\|h_j\|^2$
& $\QW^{x}/\|h_j\|^2$ & normalized difference\\
\midrule
1&0.26745365580067245&0.26745365580067245&$-1.873\cdot10^{-22}$\\
2&0.85826315683789947&0.85826315683789947&$-5.736\cdot10^{-21}$\\
3&0.28721078110996577&0.28721078110996577&$-1.506\cdot10^{-22}$\\
4&0.33249049031538647&0.33249049031538647&$-2.305\cdot10^{-22}$\\
5&0.26385733080320823&0.26385733080320823&$-3.014\cdot10^{-22}$\\
\bottomrule
\end{tabular}
\end{center}
The $c_L$ computations at 50/80-digit settings differ by
$2.12\cdot10^{-61}$ in the full run.

\section{A2: exact boundary-strip attack}

For the zero extension $E$ and
$(\tau_xf)(y)=f(y+x)$,
\[
 \langle Eh,\tau_xEh\rangle=g(x),\qquad
 \|Eh-\tau_xEh\|_{L^2(\mathbb R)}^2=2(g(0)-g(x)).
\]
Two rational step functions were used.  The first has values
$(2,-1,3)$ on the three cells cut at
$-3/10,-1/10,1/10,3/10$.  The second has values $(-2,5/2)$ on the
two half intervals.  Exact interval intersection arithmetic gives:

\begin{center}
\small
\begin{tabular}{@{}rrrrrr@{}}
\toprule
step & $x$ & $2(g_0-g_x)$ & whole-line norm
& interval-only norm & omitted strip\\
\midrule
1&$1/20$&$19/10$&$19/10$&$17/10$&$1/5$\\
1&$1/10$&$19/5$ &$19/5$ &$17/5$ &$2/5$\\
1&$1/4$ &$13/2$ &$13/2$ &$113/20$&$17/20$\\
2&$1/20$&$61/40$&$61/40$&$53/40$&$1/5$\\
2&$1/10$&$61/20$&$61/20$&$53/20$&$2/5$\\
2&$1/4$ &$61/8$ &$61/8$ &$53/8$ &$1$\\
\bottomrule
\end{tabular}
\end{center}
All six whole-line identities are exact.  Every interval-only norm is
strictly too small.  R494 uses the whole-line norm and retains the two
boundary strips.

\section{A3: the false-world test}

Replace $c_L$ by $2c_L$ while leaving $k,w,\Pi$ fixed.  This control
functional is known negative: the $h_1$ test above has Rayleigh
quotient
\[
 0.26745365580067245-c_L=-1.9249512553313268.
\]
The independent truncated-kernel chain also becomes
\[
 \kappa_w-2c_L-\lambda_{\max}(T_w-\Pi)
 =-2.188167590810956.
\]
Thus the machinery does not prove positivity in the negative control
world.  It fails by an order-one margin, as required.

\section{A4: independent numbers}

The alternative route uses Clenshaw--Curtis quadrature on the physical
square $[-L,L]^2$ and a symmetric Nystr\"om matrix.  It does not use
the r494 Gauss--Legendre/Chebyshev surrogate to obtain its diagnostic
numbers.
\[
 \kappa_w=3.69800804029897744651553\ldots .
\]

\begin{center}
\small
\begin{tabular}{@{}rrrrrr@{}}
\toprule
\CC order&$\lambda_{\max}(S)$&$\lambda_{\max}(P_{401}SP_{401})$
&$r_{401}$&$\|QTP\|_{\rm HS}$&$\|QTQ\|_{\rm HS}$\\
\midrule
513&1.501362199357172&1.501362199357172
&$5.952545\cdot10^{-4}$&$1.740576\cdot10^{-4}$&$5.419740\cdot10^{-4}$\\
769&1.501365757748623&1.501365757748621
&$6.826870\cdot10^{-4}$&$1.816422\cdot10^{-4}$&$6.325138\cdot10^{-4}$\\
1025&1.501365808845935&1.501365808845941
&$6.794869\cdot10^{-4}$&$1.804971\cdot10^{-4}$&$6.297174\cdot10^{-4}$\\
\bottomrule
\end{tabular}
\end{center}
At order 1025 the sampled Legendre Gram defect is
$6.335\cdot10^{-12}$.  The r494 section centre
$1.501365606961153$ differs by $2.02\cdot10^{-7}$ and is safely under
the outward upper bound $1.5013673873$.  The independent tail lies
below the certified r494 upper bound
$7.0457565325\cdot10^{-4}$.

An adaptive square-root endpoint transform gives the uncut
$N=36$ even diagnostic
$7.571961013293050\cdot10^{-3}$.  Combining it with the independent
cut section gives a diagnostic cutoff loss
$3.334640692256198\cdot10^{-3}$, within
$2.02\cdot10^{-7}$ of the r494 value.  The load-bearing direct
interval subtraction is
\[
 \begin{split}
 &(3.69800804029-2.19240491114-1.5013673873)\\
 &\qquad{}-3(7.0457565325\cdot10^{-4})
 =2.122014890249863\cdot10^{-3}>2.1\cdot10^{-3}.
 \end{split}
\]

\section{A5: the full operator-class check}

Let $P=P_{401}$ and $Q=I-P$.  Because $w=0$ near zero and the quintic
ramp matches through two derivatives, $w(x)k(x)$ is $C^2$ on
$[0,2L]$.  Hence
\[
 T_w(y,z)=\tfrac12w(|y-z|)k(|y-z|)
\]
is bounded on a compact square and is Hilbert--Schmidt.  Smoothness
supplies the observed rate; Hilbert--Schmidt membership alone supplies
convergence.

Normalized Legendre polynomials are a complete orthonormal basis of
$L^2[-L,L]$.  Parseval in the Hilbert--Schmidt class therefore gives
the exact identity
\[
 r_n^2=\|T_w\|_{\rm HS}^2-\sum_{i,j<n}|(T_w)_{ij}|^2
 =2\|Q T_wP\|_{\rm HS}^2+\|Q T_wQ\|_{\rm HS}^2.
\]
Thus the old off-space problem cannot hide in an unrepresented class:
the sum contains \emph{both} off-blocks and the entire block-diagonal
rest beyond $n$.  Since operator norm is at most Hilbert--Schmidt norm,
\[
 2\|PT_wQ\|+\|QT_wQ\|\le3r_n.
\]
The independent order-1025 decomposition finds the block-diagonal
rest explicitly:
\[
 \|QT_wP\|_{\rm HS}=1.8049712\cdot10^{-4},\qquad
 \|QT_wQ\|_{\rm HS}=6.2971739\cdot10^{-4}.
\]
It is not small enough to discard, but it is included.  Directly,
\[
 2\|QT_wP\|_{\rm HS}+\|QT_wQ\|_{\rm HS}
 =9.9071162\cdot10^{-4}
 <3(7.0457565325\cdot10^{-4}).
\]
The rank-two pole vectors are entire.  On $|x/2|\le0.15$, the
degree-400 Taylor remainder is bounded by
$e^{0.15}0.15^{401}/401!$, far below $10^{-30}$ even after the
rank-two operator conversion.  No pole remainder is hidden.

\section{Kill table and verdict}

\begin{center}
\small
\begin{tabular}{@{}p{0.08\linewidth}p{0.79\linewidth}@{}}
\toprule
Attack&Result\\
\midrule
A1&Survives: kernel, $c_L$, $g(0)$ term, pole sign, and five direct
frequency/$x$ comparisons agree.\\
A2&Survives: the zero-extension identity is exact; both boundary-strip
masses are retained.\\
A3&Survives: a known negative doubled-$c_L$ world makes both the
Rayleigh quotient and the certificate budget negative.\\
A4&Survives: independent quadrature reproduces $\kappa_w$, section
eigenvalue, cutoff loss, tail, and final arithmetic inside all outward
bounds.\\
A5&Survives: completeness plus Hilbert--Schmidt Parseval includes both
off-blocks and the nonzero $QTQ$ remainder.\\
\bottomrule
\end{tabular}
\end{center}

\begin{theorem}[Red-team verdict]
The r494 assertion survives all five attacks:
\[
 \boxed{\QW(h)\ge2.1\cdot10^{-3}\|h\|_2^2
 \quad\text{for }\operatorname{supp}h\subset[-0.3,0.3].}
\]
The red-team verdict is
\texttt{CONFIRMED($c=2.1000\mathrm e{-3}$)}.
\end{theorem}

No RH claim and no anti-RH claim is made.

\end{document}
